\section{Exercices}\label{exercices} % attendu: C5

\subsection{Premier}\label{sec:premier} % attendu: C5

\begin{exercise}[label=exo:cauchy] % attendu: C5
  Montrer qu'une suite de Cauchy est bornée.
\end{exercise}

\begin{example}[label=ex:harmonique] % attendu: C5
  La série harmonique diverge.
\end{example}

\begin{myremark}[label=rmk:alias] % attendu: C5
  Un alias du template suit sa cible : `myremark` est un `remark`, donc `rem:`.
\end{myremark}

\begin{align}
  S_n &= \sum_{k=1}^n \frac1k \label{eqn:somme} % attendu: C5
\end{align}

Voir la section~\ref{exercices}, \ref{sec:premier}, l'Exercice~\ref{exo:cauchy},
l'Exemple~\ref{ex:harmonique}, la Remarque~\ref{rmk:alias} et \eqref{eqn:somme}.
